\begin{definition}[Exchanging actual internal reference registers]
    \label{def:peps_regular_charge_move_coordinates}
    \lean{TNLean.PEPS.regularChargeMoveCoordinatePermutation,
        TNLean.PEPS.regularChargeMoveParameter}
    \leanok
    \uses{def:peps_regular_coordinate_physical_permutation}
    Let $R$ be a connected finite block, with a spanning tree and root.
    For two internal bonds $e,f$, exchange $a_e$ and $a_f$ in its actual
    accessible coordinates $(y,a,r,z)$, retaining all other registers.
    For actual group operators $u$, let $k$ be the derived rooted tree gauge.
    The corresponding transported parameter is
    $p'=p k_{\mathrm{tail}(e)}k_{\mathrm{tail}(f)}^{-1}$.
    This is the register exchange used for charge motion in
    \cite{Schuch2010PEPS}, lines 2489--2507.
\end{definition}

\begin{theorem}[One original-spin unitary moves a literal charge]
    \label{thm:peps_regular_physical_charge_motion}
    \lean{TNLean.PEPS.regularChargeMoveCoordinatePermutation_coordinateTranslation,
        TNLean.PEPS.exists_unitary_regularPhysicalChargeMotion}
    \leanok
    \uses{def:peps_regular_charge_move_coordinates,
        thm:peps_regular_coordinate_permutation_support,
        thm:peps_regular_weighted_open_recovery,
        thm:peps_regular_physical_unitary_transport}
    Suppose the original sites are locally $G$-isometric for the regular
    incident-edge action. One unitary $W_{e,f}$ on the original spins of $R$
    is chosen before every actual group-valued background $u$, coefficient
    function $\chi$, parameter $p$ and boundary configuration $\theta$.
    It satisfies
    \[
        W_{e,f}\Psi_R(u;\chi,p,e,\theta)
        =\Psi_R(u;\chi,p',f,\theta),
    \]
    where $\Psi_R$ is the literal weighted open contraction with weight
    $\chi(p\eta_e)$, and $p'$ is the derived tree transport above.
    The character is retained when $\chi$ is an irreducible character.
    No coefficient identity or Gram matrix is assumed.
\end{theorem}
\begin{proof}\leanok
    On surviving canonical terms the charge label is
    $k_{\mathrm{tail}(e)}x^{-1}a_e$. After the swap its target value is
    $k_{\mathrm{tail}(f)}x^{-1}a_e$; the parameter $p'$ cancels the
    intervening tree gauge. The swap commutes with all independent vertex
    translations and their average. Local isometry therefore lifts it to
    a unitary on the original spins, and the actual weighted recovery
    identity gives the stated contraction action.
\end{proof}

\begin{theorem}[The weighted charge is the literal tensor insertion]
    \label{thm:peps_regular_weighted_charge_contraction}
    \lean{TNLean.PEPS.regularWeightedOpenRegionWeight_eq_openRegionWeight_regularEdgeCharacterSite}
    \leanok
    \uses{thm:peps_regular_weighted_open_recovery}
    For an internal bond $e$ and identity group operators, the weighted open
    contraction with weight $\chi(p\eta_e)$ equals the ordinary open tensor
    contraction with the diagonal insertion $\operatorname{diag}\chi(p\cdot)$
    on the ordered tail of $e$. Every boundary label is retained.
    This is the charge insertion of \cite{Schuch2010PEPS}, lines
    2432--2453.
\end{theorem}
\begin{proof}\leanok
    The identity bond shares one group label at both endpoints. The product
    of the site factors differs from the original product by precisely the
    one factor $\chi(p\eta_e)$ at its tail. Summing over the actual incident
    edge labels gives the claimed identity, with the same boundary constraint.
\end{proof}

\begin{theorem}[Untwisted motion without supplied tree choices]
    \label{thm:peps_regular_untwisted_physical_charge_motion}
    \lean{TNLean.PEPS.regularRegionTreeGauge_one,
        TNLean.PEPS.exists_unitary_regularPhysicalChargeMotion_of_connected}
    \leanok
    \uses{thm:peps_regular_physical_charge_motion,
        thm:peps_regular_weighted_charge_contraction}
    If the group operators are the identity, the actual rooted gauge is
    the identity. For any connected finite block and two internal bonds,
    one original-spin unitary therefore moves the literal diagonal insertion
    from $e$ to $f$ in the ordinary open tensor contraction, retaining $\chi,p$ and every boundary configuration. The tree and root
    are chosen in the proof. This is a finite-block charge-motion consequence
    of \cite{Schuch2010PEPS}, lines 2489--2507.
\end{theorem}
\begin{proof}\leanok
    Choose a spanning tree from the actual induced connectivity. Its identity
    gradient has the identity normalized vertex labels. Thus $p'=p$ in the
    preceding physical theorem.
\end{proof}
